\documentclass[12pt]{article}
\usepackage[utf8]{inputenc}
\usepackage[T1]{fontenc}
\usepackage{hyperref}
\title{Paraguayan Sign Language Language Resources for K-12 Curriculum Development}
\author{Compiled by Claude Sonnet 5.5 (claude-sonnet-5-5) for Andr\'as Kornai}
\date{}
\begin{document}
\maketitle

\section{Introduction}
Paraguayan Sign Language (Lengua de Se\~nas Paraguaya, LSPy; ISO 639-3 pys; Glottocode para1318) is the sign language of the Paraguayan deaf community. Search results indicate that Law 6530/20 gives it official recognition as a language of communication and instruction. The search found little in the way of ready-made K-12 curriculum materials. What exists is mostly news coverage, a digital sign dictionary and some picture-book videos. This report is short because little else turned up.

\section{Primary Educational Resources}
\begin{itemize}
\item UNICEF Paraguay story on classic children's literature for deaf and hearing children, with LSPy video books produced with the Ministry of Education: \url{https://www.unicef.org/paraguay/stories/m%C3%A1s-cl%C3%A1sicos-de-la-literatura-para-ni%C3%B1os-sordos-y-oyentes}
\item Report on a school for deaf students and the importance of sign language (ABC Color, 2023): \url{https://www.abc.com.py/nacionales/2023/09/22/escuela-de-sordos-comparte-importancia-de-la-lengua-de-senas/}
\end{itemize}

\section{Linguistic and Reference Materials}
\begin{itemize}
\item News item on a planned digital sign-language dictionary (Hoy): \url{https://www.hoy.com.py/nacionales/lengua-de-senas-tendra-diccionario-digital}
\item Academic article on a trilingual (Spanish, Guarani, LSPy) literacy proposal for deaf people (Universidad Espa\~nola de Paraguay yearbook): \url{https://ojs.uep.edu.py/index.php/anuarioacademico/article/download/67/50}
\end{itemize}

\section{Religious and Community Resources}
\begin{itemize}
\item Report on Christian churches joining to translate the Bible into sign language (La Naci\'on, 2026): \url{https://www.lanacion.com.py/pais/2026/09/30/iglesias-cristianas-se-unen-para-traducir-la-biblia-a-lengua-de-senas/}
\item Campaign promoting sign language (La Naci\'on, 2018): \url{https://www.lanacion.com.py/pais/2018/09/14/presentan-campana-que-promociona-lengua-de-senas/}
\end{itemize}

\section{Digital and Audio Resources}
\begin{itemize}
\item Basic sign language primer article (La Naci\'on, 2019): \url{https://www.lanacion.com.py/pais/2019/09/23/lenguaje-de-senas-lo-basico-que-todos-debemos-saber/}
\item Article on including about 20,000 deaf people (\'Ultima Hora): \url{https://www.ultimahora.com/manos-que-escriben-el-aire-incluir-20000-personas-sordas-n732324}
\end{itemize}

\section{Curriculum Development Context}
Search results suggest that few deaf children reach formal early education, that interpreters are scarce, and that no schools teach LSPy to hearing learners. A K-12 curriculum would have to start from the deaf community's own organisations and educators. The survey found no complete textbooks or curriculum documents.
\end{document}